\chapter*{Abstract}
\addcontentsline{toc}{chapter}{Abstract}
\thispagestyle{plain}

A mobile phone is always attached to one antenna, and the operator writes a line
down each time it changes. This report asks whether a language model, retrained
on those anonymous lines, can guess which antenna a person's phone will be
attached to next, better than the rule ``it is still on the same one as a moment
ago''. That rule, called the \emph{persistence baseline} here, is the yardstick
for every claim in this document, and it is a hard one: on an ordinary day of
records it is right 69.2 times out of a hundred, and for five of the eleven
iterations of this internship nothing I trained beat it.

The material is fifteen daily exports covering the Karlovy Vary study area, 369
antennas mounted on 113 masts, with one day of records per person. On the
thousand people held out of the largest set, the retrained model answers
correctly \textbf{4.3 more times out of a hundred} than the one-line rule, 54.5
against 50.2. What produced that gain was neither a bigger model nor more data,
but conditioning on \emph{who the person is}. Two further results frame it:
population size stops paying at around two thousand people, and 22.4 correct
answers per hundred separate the stay-at-home user from the morning-moving user,
against 0.4 across the entire range of training sizes. Difficulty lives in the
person, not in the dataset.

\vspace{0.8em}
\noindent\textbf{Keywords:} human mobility prediction, call detail
records,\footnote{\textbf{CDR}, \emph{call detail record}: the line an operator
writes whenever a phone and the network interact. It holds an anonymous
identifier, a time, and the antenna in use. No content of any kind.} next-cell
prediction, large language models, LoRA, QLoRA, behavioural clustering, scaling.
